\documentclass[a5paper,10pt]{article}

% https://github.com/InDevRus/filippov-solutions

% Нумерация страниц
\pagenumbering{gobble}

% Настройка полей
\usepackage{geometry}
\geometry{left=1cm}
\geometry{right=1cm}
\geometry{top=1cm}
\geometry{bottom=1.5cm}

% Вставка таблицы
\usepackage{array}
\usepackage{tabularx}

% Инструменты выравнивания формул
\usepackage{amsmath}
\usepackage{mathtools}

% Диакритика в математике
\usepackage{amsfonts}

% Вычисления размеров на ходу
\usepackage{calc}

% Удобные записи производных
\usepackage[italic=true]{derivative}

% Настройки сносок
\usepackage{enotez}

% Длинные знаки векторов
\usepackage{anyfontsize}
\usepackage[b]{esvect}

% Вставка картинок
\usepackage{graphicx}

% Вставка красной строки
\usepackage{indentfirst}
\setlength{\parindent}{1em}

% Условия в командах
\usepackage{xifthen}

% Луа-скрипты
\usepackage{luacode}

% Настройка команд отдельно в математическом и текстовом режимах
\usepackage{mathcommand}

% Для повторения LaTeX кода
\usepackage{multido}

% Конфигурация отступов формул
\delimitershortfall-1sp
\usepackage{mleftright}
\mleftright

% Растягивание объектов
\usepackage{relsize}

% Обтекание текстом
\usepackage{wrapfig}
\usepackage{subcaption}
\usepackage{floatflt}

% Поддержка ввода кириллицы
\usepackage[english, russian]{babel}

% Настройка корневой позиции для компилятора
\usepackage{import}
\providecommand*{\ProjectRootPath}{..}

\import{\ProjectRootPath/preambles/macros/}{theorems}

\import{\ProjectRootPath/preambles/fonts}{font_setup}

% Знак градуса
\usepackage{siunitx}

\import{\ProjectRootPath/preambles/macros/}{operators}
\import{\ProjectRootPath/preambles/macros/}{redefinitions}

\import{\ProjectRootPath/preambles/fonts}{letters_setup}
\import{\ProjectRootPath/preambles/fonts/adjustments}{custom_superscripts}

\import{\ProjectRootPath/preambles/custom}{formatting}
\import{\ProjectRootPath/preambles/custom}{frames}
\import{\ProjectRootPath/preambles/custom}{list_setup}

% TODO: перевести команды на современный стиль объявления

\begin{document}

\begin{framed}
    Для угла $\phi = \ang{90}$ имеем 
    $\tg{\beta} = \tg\br{\alpha \pm \phi} 
    = \tg\br{\alpha \pm \ang{90}} = -\ctg \alpha$.
\end{framed}

$\pow{x}{2} = y + Cx$.
$C = \dfrac {\pow{x}{2} - y} {x}$.
$0 = \dfrac {\br{2x - y\`}x - \br{\pow{x}{2} - y}} {\pow{x}{2}} 
= \dfrac {\pow{x}{2} - xy\` + y} {\pow{x}{2}}$.
Имеем дифференциальное уравнение: 
$y\` = \dfrac {y + \pow{x}{2}} {x}$.
$\tg\alpha = \dfrac {y + \pow{x}{2}} {x}$.
$\tg\beta = -\ctg\alpha = -\dfrac {x} {\pow{x}{2} + y}$.
Окончательно, уравнение изогональных (ортогональных) 
траекторий имеет вид $z\` = -\dfrac {x} {\pow{x}{2} + z}$.

\begin{remark}
    Уравнение $z\` = -\dfrac {x} {\pow{x}{2} + z}$ имеет общее 
    решение, например, в неявном виде: 
    $\pow{x}{2} = \dfrac {1} {2} - z + \widetilde{C} \pow{e}{-2z}$,
    $\widetilde{C} > -\dfrac {1} {2}$.
\end{remark}

\end{document}
